\documentclass[10pt,a4paper]{amsart}
\usepackage[margin=19mm]{geometry}
\usepackage{amsmath,amssymb,mathtools}
\usepackage[scr=rsfs,cal=euler]{mathalfa}
\usepackage{xfrac}
\usepackage[unicode,hidelinks,bookmarks=false]{hyperref}
\newcommand{\ie}{i.e.\ }
\newcommand{\cf}{cf.\ }
\newcommand{\resp}{respectively\ }
\newcommand{\Subsection}{\subsection}
\providecommand{\nobreakdash}{\nobreak\hbox{-}}
\setlength{\emergencystretch}{2em}
\multlinegap0pt
\usepackage{bm}
\usepackage{bbm}
\usepackage{dsfont}
\usepackage{xspace}
\usepackage{cancel}
\usepackage{mathtools}
\usepackage[shortlabels]{enumitem}
\usepackage{amsthm}
\makeatletter
\@ifundefined{thm}{\newtheorem{thm}{Theorem}}{}
\@ifundefined{lem}{\newtheorem{lem}[thm]{Lemma}}{}
\makeatother
\makeatletter
\newcommand{\Haus}{{\text{\rm Haus}}}
\newcommand{\LL}{\mathcal{L}}
\newcommand{\MM}{\mathcal{M}}
\newcommand{\bB}{\mathbf{B}}
\newcommand{\bQ}{\mathbf{Q}}
\newcommand{\bT}{\mathbf{T}}
\newcommand{\bbF}{\mathbb{F}}
\newcommand{\bbZ}{\mathbb{Z}}
\newcommand{\cP}{\mathcal{P}}
\expandafter\ifx\csname claimproof\endcsname\relax
\newenvironment{claimproof}[1][\claimproofname]{\begin{proof}[#1]\renewcommand*{\qedsymbol}{\(\bullet\)}}{\end{proof}}
\fi
\newcommand{\eps}{\varepsilon}
\newcommand{\hto}{\stackrel{\Haus}{\longrightarrow}}
\newcommand{\leqnomode}{\tagsleft@true\let\veqno\@@leqno}
\newcommand{\linkdest}[1]{\Hy@raisedlink{\hypertarget{#1}{}}}
\makeatother
\title[Convergent sequences of combinatorial submodular setfunction]{Convergent sequences of combinatorial submodular setfunctions}
\author{Krist\'of B\'erczi, M\'arton Borb\'enyi, L\'aszl\'o Lov\'asz, L\'aszl\'o M\'arton T\'oth}
\date{Source: arXiv:2507.15105v1 (CC BY 4.0)}
\begin{document}
\maketitle

\noindent\emph{Source and licence.} Convergent sequences of combinatorial submodular setfunctions. Version arXiv:2507.15105v1 (CC BY 4.0), distributed under the Creative Commons Attribution 4.0 International licence, as recorded in the arXiv licence metadata for this exact version and shown on the abstract page https://arxiv.org/abs/2507.15105v1. Full rights notice: licenses/arxiv-2507.15105-CC-BY-4.0.txt.\par\smallskip

\noindent\textit{Editorial scope.} Counted unit: the Thm whose source label is \texttt{THM:FIN-GEOM} in the article (source0.tex lines 618--622, source label \texttt{THM:FIN-GEOM}), delivered together with the article's own complete original proof of it (source0.tex lines 623--657, one \texttt{proof} environment of 35 lines). No mathematics is written, repaired or re-derived: statement and proof are the article's own text line for line. Required context retained uncounted: THM:MAT-CONV (lines 449--459); eq:rkm (lines 527--530); LEM:M-A-BASE2 (lines 564--574). Every definition, display and label of those lines is retained unchanged. Omitted from this package and not redistributed: 1--448, 460--526, 531--563, 575--617, 658--1017. Nothing inside a retained excerpt is omitted or altered: every retained body is delivered exactly as the article's lines are written, so no omission receipt of any kind accompanies this package. Citations: the retained excerpts carry no citation command at all, so no bibliography entry of the article is transcribed and no reference is redistributed. Reference adaptation: none was needed and none was made; every reference inside the delivered unit resolves inside it, so no unresolved reference or citation is produced. Printed slips kept literal: no printed word, symbol, hypothesis or number of the counted statement or of its proof was altered. Layout and class: the article's own class is \texttt{article} with packages including fullpage, inputenc, amsthm, amssymb, amsmath, bbm, soul, caption, subcaption, epsfig, graphicx, graphics, color, tikz; the wrapper is the lane's pinned \texttt{amsart} editorial wrapper at 10pt on A4, which restates the theorem-like environments the retained text needs and adds the article's own macro definitions that the retained text needs (\texttt{\textbackslash Haus}, \texttt{\textbackslash LL}, \texttt{\textbackslash MM}, \texttt{\textbackslash bB}, \texttt{\textbackslash bQ}, \texttt{\textbackslash bT}, \texttt{\textbackslash bbF}, \texttt{\textbackslash bbZ}, \texttt{\textbackslash cP}, \texttt{\textbackslash eps}, \texttt{\textbackslash hto}, \texttt{\textbackslash leqnomode}, \texttt{\textbackslash linkdest}), with extra wrapper packages (bm, bbm, dsfont, xspace, cancel, mathtools, enumitem). The delivered numbering is the unit's own. Assets: no retained span contains a figure environment, a graphics-inclusion command, a PDF-page inclusion, a table or a TikZ picture; no image file is copied into this package, the package declares no asset dependency and no image file is redistributed with it. Licence: version arXiv:2507.15105v1 of this article is distributed under the Creative Commons Attribution 4.0 International licence, as recorded in the arXiv licence metadata for this exact version and shown on the abstract page https://arxiv.org/abs/2507.15105v1. A full rights notice is delivered at licenses/arxiv-2507.15105-CC-BY-4.0.txt. Characters: every retained excerpt is pure ASCII, so the delivered engine, XeLaTeX with its default Latin Modern text font, reports no missing character.
\begin{thm}\label{THM:MAT-CONV}
Let $(\MM_n=(E_n,r_n)\colon n\in\bbZ_+$) be a sequence of matroids with $r_n(E_n)=n$, and
normalized rank functions $\rho_n=r_n/n$. Let $\LL_n$ be the family of flats in
$\MM_n$. Assume that
\begin{enumerate}[label=(\roman*)]\itemsep0em
    \item for all $m\le n$, there exists a rank preserving embedding $\LL_m\to\LL_n$, and \label{it:i}
    \item if $m\mid n$, then there exists a stretch embedding $\LL_m\to\LL_n$. \label{it:ii}
\end{enumerate}
Then $(\rho_n\colon n\in\bbZ_+)$ is crop-convergent, and  the limit of
$k$-crops is given by $T_k(\rho_n)\hto \bT_k=\bigcup_{n\in\bbZ_+} T_k(\rho_n)$ for all $k\in\bbZ_+$.
\end{thm}

\begin{equation*}\leqnomode
    \text{For any two flats $F\subseteq A$ where $r(A)\ge m$, we have
$|A\setminus F|\ge k\cdot (r(A)-r(F))$.} \tag*{(R(k,m))}\linkdest{rkm}{} \label{eq:rkm}
\end{equation*}

\begin{lem}\label{LEM:M-A-BASE2}
Let $\MM=(E,r)$ be a matroid satisfying~\ref{eq:rkm} for some $m\ge k\ge 1$.
Then
\[
d^\Haus\bigl(T_k(\rho),T_k^\Delta(\rho)\bigr)\le \frac{k\cdot m}{r(E)}
\]
and
\[
d^\Haus\bigl(T_k^\nabla(\rho),Q_k(\rho)\bigr)\le \frac{k\cdot m}{r(E)}.
\]
\end{lem}

\begin{thm}\label{THM:FIN-GEOM}
Let $ \bbF$ be a finite field, $|\bbF|=q$, and let $\rho_n=\dim/n$ denote the
normalized rank function of the linear space $\bbF^n$. Then the sequence
$(\rho_n\colon n\in\bbZ_+)$ is quotient-convergent.
\end{thm}

\begin{proof}
We start by observing that the sequence $(\bbF^n\colon n\in\bbZ_+)$ has two important properties. First, if $m\le n$, then the map $\LL_m\to\LL_n$ defined by $A\mapsto A\oplus \mathbf{0}$, where $\mathbf{0}$ is the zero subspace of $\bbF^{n-m}$, is a rank-preserving lattice embedding. Second, if $m\mid n$, then the map $\LL_m\to\LL_n$ defined by $A\mapsto A\oplus\dots\oplus A$, with $n/m$ terms, is a stretch embedding. 

It follows that the sequence $(\bbF^n\colon n\in\bbZ_+)$ satisfies the conditions of Theorem~\ref{THM:MAT-CONV}, and the sequence $(T_k(\bbF^n)\colon n\in\bbZ_+)$ is
convergent in the Hausdorff distance; let $ \bT_k$ be its limit. It is not hard to verify that every matroid $\bbF^n$ satisfies condition
\hyperlink{rkm}{($R(k,2k)$)} for every $k$ and $n$, and hence Lemma~\ref{LEM:M-A-BASE2} implies
that the sequence $(T^\Delta_k(\bbF^n)\colon n\in\bbZ_+)$ converges to $ \bT_k$ in the
Hausdorff distance as well. We claim that the sets $Q_k(\rho_n)$ converge in Hausdorff distance to the set
\[
 \bQ_k=\{\psi\in \bT_k\colon \max_i\psi(\{i\})=1\}
\]
for every $k\in\bbZ_+$. Let $\psi\in Q_k(\rho_n)$, and write $\psi=\rho_n/\cP$ where $\cP=(A_1,\dots,A_k)$ is a measurable partition of $ \bbF^n$. Clearly,
there exists $i\in[k]$ such that $|A_i|\ge q^n/k$. This implies that
$\dim(A_i)\ge n-\log_q k$ and hence
\begin{equation*}\label{EQ:QQK1}
\max_{i\in[k]} \psi(\{i\}) \ge  1-\frac{\log_q k}{n}.
\end{equation*}
This, together with $Q_k(\rho_n)\subseteq T_k(\rho_n)$, implies that every limit point of the sequence
$Q_k(\rho_n)$ is in $\bQ_k$. On the other hand, let
$\psi\in \bQ_k$, where we may assume without loss of generality that $\psi(\{k\})=1$. This implies that $\psi(X)=1$ for every $X\subseteq [k]$ such that $k\in X$.

Let $\eps>0$. We know that $\psi \in \bT_k$, so there exists a family
$ \bB=(B_1,\dots,B_k)$ of disjoint subsets of $\bbF^n$ such that
$\|\rho_n/ \bB-\psi\|\le \eps$. Let us replace $B_k$ by the set
$B_k'=\bbF^n\setminus\bigcup_{j=1}^{k-1} B_j$, then we get a partition $\cP$ of
$\bbF^n$ such that $\rho_n/\cP\ge\rho_n/ \bB$, and so
\[
\|\rho_n/\cP-\rho_n/ \bB\|=\rho_n(B'_k)-\rho_n(B_k) \le \frac{\log_q k}{n}.
\]
It follows that 
\[
\|\rho_n/\cP-\psi\| \le \eps+\frac{\log_q k}{n} =\eps+o(1),
\]
proving that $Q_k(\bbF^n)\hto \bT_k$ as $n\to\infty$.
\end{proof}

\end{document}
